\chapter{结论与展望}
\label{ch:conclusion}

\section{全文工作总结}

本文围绕资源受限无人机网络中的协同通信问题，研究了能量感知路由、选择性语义通信和不完备反馈下的可靠交付方法。无人机网络一方面需要在移动拓扑和有限能量下维持数据回传，另一方面需要在有限通信机会中选择有助于任务判断的局部证据；当这些证据进入实际传输过程后，还需要处理发送失败与确认信息不完整造成的接收状态不确定。本文沿上述过程逐步细化决策对象，分别建立模型、设计算法并开展实验分析。

在能量感知路由方面，本文针对关键中继容易承受过多转发任务的问题，建立了包含动态连接、数据队列和能量支出的飞行自组网模型。通过李雅普诺夫漂移加惩罚方法，把长期能量约束转化为逐时隙可使用的资源压力，并据此构造 LEADR 分布式路由与功率决策。算法根据候选链路速率和能量代价计算发送功率，再结合剩余能量及地理进度选择下一跳，使转发任务能够随网络状态变化进行调整。理论分析说明了功率子问题的求解条件和虚拟队列与累计预算的联系，实验从节点数量、通信范围和飞行速度考察网络表现。该工作表明，路由评价需要同时观察转发效率、丢包、能量分布和服务周期，才能判断局部传输优势是否能够持续。

在选择性语义通信方面，本文针对不同局部观测和消息粒度具有不同任务作用的问题，提出 EviComRL。通过冻结视觉语言模型构造任务相关证据，将语义提取过程与轻量通信策略学习分离；在候选特征中加入中心信念、链路状态、可靠性和已选历史，使策略能够根据当前任务状态逐步调整消息选择。算法以动作掩码表示预算与接收约束，结合代理预热和 PPO 学习通信与停止行为。主实验设置中，方法取得 58.17\% 的准确率、50.58\% 的宏平均 F1 和 49.70 token 的每任务平均通信 token 数。跨条件结果说明，观测覆盖决定可获得证据的基础，消息粒度和预算影响质量成本取舍，链路恶化可能在分类标签变化较小时明显降低有效信念质量。可靠性特征的重复实验进一步支持利用可观测信念分布和节点共识辅助识别证据来源。

在语义可靠交付方面，本文把研究推进到给定候选之后的发送过程，提出 C-BFC。针对同一分类结果可能被学习模型赋予不同效用的问题，构造同类别一致性投影；针对未确认发送不能直接说明数据失败的问题，利用累计位图和发送历史递推实际接收集合的概率分布。在此基础上，通过两步发送前瞻和等待反馈的信息价值选择新发、重传、等待或停止，并使用虚拟能量队列与剩余额度约束跨任务支出。在十七个主比较条件下，方法的平均按时正确完成率为 49.84\%，单位正确任务能耗为 3.615 mJ。与 SD-SAC 和 TP-DQN 相比，能耗分别降低 16.59\% 和 14.32\%，同时保持略高的平均正确完成率。消融结果说明，接收信念有助于减少对确认缺失的误判，等待用部分时间换取更合适的后续动作，一致性修正主要改善低收益通信支出。

三项工作的共同认识是，通信资源的价值取决于其所支持的后续过程。高质量链路需要与节点能量持续性协调，语义内容需要与接收端当前需求协调，重传和等待需要与接收状态及剩余任务机会协调。本文采用队列优化、策略学习和显式信念控制处理不同形式的不确定性，形成了从网络传输到任务交付的研究体系。各方法的实验结果在相应模型和评价范围内成立，为后续开展完整无人机协同系统的联合设计提供了基础。

\section{后续研究展望}

\subsection{移动链路与语义控制的进一步耦合}

本文的三个研究章分别在与其决策对象匹配的环境中开展评价，完整系统还需要在路由状态、语义候选和交付反馈之间建立持续运行的接口。后续可把路由产生的排队时延、有效速率和节点存活状态直接用于语义候选的代价预测，使上层选择能够响应真实路径变化；也可以把任务的重要程度和剩余截止时间反馈给路由层，使底层服务优先级与语义需求保持一致。该方向需要处理路由时隙、语义任务与反馈事件之间的时间尺度差异，并量化状态交换本身的通信开销。

\subsection{更丰富的语义任务与接收状态表示}

本文采用固定类别空间的证据可见性任务，使内容选择与交付策略的作用能够在共同语义基础上比较。后续可以扩展到目标定位、地图更新、异常事件判断和生成式回答等任务，使用与实际损失相对应的语义效用。在连续估计或自然语言输出中，不同证据集合之间的等价关系不再只由预测类别决定，需要重新设计效用一致性约束。同时，当候选数量增加时，显式接收集合的规模呈指数增长，可研究基于候选聚类、粒子表示或分层筛选的近似方法，在状态表达能力与在线计算成本之间取得平衡。

\subsection{相关信道、反馈竞争与质量要求}

实际无线环境中的数据丢失可能具有时间相关性，多个节点的反馈也可能发生竞争或碰撞。后续可使用测量轨迹描述相关衰落和共享控制信道，进一步研究确认缺失携带的统计信息，以及反馈频率、反馈长度与传输调度之间的关系。对于具有明确最低服务质量的任务，可将当前加权收益目标扩展为带质量约束的资源分配问题，分析不同质量要求下能耗、时延和连续完成能力的可行范围。该研究需要在充分验证约束可行性的基础上确定控制参数，避免仅通过提高惩罚权重近似严格需求。

\subsection{机载部署与完整资源核算}

后续系统实现还需要同时测量语义提取、候选整理、策略推理、数据发送和反馈处理的实际开销。冻结视觉语言模型能够减少策略训练阶段的重复调用，但机载运行仍受到模型大小、推理时延和功耗的约束。可进一步结合模型压缩、边缘协同推理和计算通信联合调度，观察轻量控制器带来的通信节约能否覆盖额外计算成本。在真实设备或硬件在环环境中复现任务过程，也有助于识别队列同步、反馈时效和资源状态估计对算法效果的影响，从而把本文的方法推进为可持续运行的无人机协同通信机制。
